\section{引言：为什么要理解 GPU}

本节课有两个核心学习目标：（1）理解 GPU 的工作原理，让 CUDA 和 GPU 不再``神秘''；（2）掌握加速算法的基本思路，能够针对新架构设计高性能实现。

\begin{importantbox}{课程两大目标}
\begin{enumerate}
    \item 理解 GPU 何时变慢、为什么变慢——揭开矩阵乘法性能曲线中那些``诡异波动''的谜底
    \item 掌握加速算法设计的核心思路——理解 FlashAttention 等算法为何能大幅提升性能
\end{enumerate}
\end{importantbox}

\begin{figure}[H]
\centering
\includegraphics[width=0.95\textwidth]{slides-images/slide-001.jpg}
\caption{课程大纲与目标：左侧为方阵乘法的 TF/s 性能曲线（标注了 Tiling、Compute Intensity、Wave Quantization 三大现象），右侧为 FlashAttention 的 tiled 计算示意图及性能对比\protect\footnotemark}
\end{figure}
\footnotetext{Slides 第2页。矩阵乘法性能数据来自 \url{https://www.thonking.ai/p/what-shapes-do-matrix-multiplications}；FlashAttention 图来自 Dao et al.}

课程的整体组织分为三部分：

\begin{enumerate}
    \item \textbf{Part 1}：GPU 深度剖析——硬件结构与关键组件
    \item \textbf{Part 2}：理解 GPU 性能——如何让 ML 工作负载跑得更快
    \item \textbf{Part 3}：综合应用——拆解 FlashAttention
\end{enumerate}

\subsection{计算驱动性能提升}

Scaling laws 告诉我们：更多的计算资源能带来可预测的语言模型性能提升。Kaplan 等人的研究表明，验证损失与计算量之间存在幂律关系 $L = 2.57 \cdot C^{-0.048}$。

\begin{figure}[H]
\centering
\includegraphics[width=0.95\textwidth]{slides-images/slide-004.jpg}
\caption{Neural Scaling Laws：计算量（PetaFLOP/s-days）与验证损失的关系\protect\footnotemark}
\end{figure}
\footnotetext{Slides 第5页。图来自 Kaplan et al., Neural Scaling Laws。}

更快的硬件、更好的利用率、更高效的并行化，这些才是真正驱动性能进步的因素。因此，理解硬件是构建高效语言模型的必要前提。

\subsection{从 Dennard Scaling 到并行计算}

在半导体发展的早期（1980s--2000s），芯片依靠 \textbf{Dennard Scaling} 来提升性能：更小的晶体管 $\rightarrow$ 更高的时钟频率 $\rightarrow$ 更低的功耗 $\rightarrow$ 更高的性能。但 Dennard Scaling 在 2000 年代初期逐渐失效，单线程性能增长趋于平缓。

\begin{figure}[H]
\centering
\includegraphics[width=0.95\textwidth]{slides-images/slide-005.jpg}
\caption{Hennessy \& Patterson 经典图表：42年处理器数据，单线程性能（蓝色）约2000年后趋平\protect\footnotemark}
\end{figure}
\footnotetext{Slides 第6页。}

\begin{knowledgebox}{Dennard Scaling 与 Moore's Law}
\textbf{Moore's Law}：晶体管数量每1--2年翻倍（仍在延续）。\\
\textbf{Dennard Scaling}：更小的晶体管可以在更高频率下以更低功耗运行（已失效）。\\
失效的原因：漏电流增大、散热极限。此后，性能提升转向\textbf{并行化}。
\end{knowledgebox}

GPU 并行计算的扩展在过去 10 年实现了超过 \textbf{1000x} 的性能提升（从 K20X 到 H100）。这一增长来自四个维度：

\begin{itemize}
    \item \textbf{数值表示}（FP32 $\rightarrow$ FP16 $\rightarrow$ Int8）：$\sim$16x
    \item \textbf{复杂指令}（DP4, HMMA, IMMA）：$\sim$12.5x
    \item \textbf{制程进步}（28nm $\rightarrow$ 5nm）：$\sim$2.5x
    \item \textbf{稀疏性}：$\sim$2x
\end{itemize}

\begin{figure}[H]
\centering
\includegraphics[width=0.95\textwidth]{slides-images/slide-006.jpg}
\caption{Bill Dally HotChips 演讲：单芯片推理性能 10 年增长 1000x\protect\footnotemark}
\end{figure}
\footnotetext{Slides 第7页。}

\begin{importantbox}{没有 GPU 扩展就没有 LLM 扩展}
``There is no LLM scaling without GPU scaling.'' GPU 的并行计算能力是大语言模型得以存在的硬件基础。
\end{importantbox}

\subsection{本章小结}

半导体性能提升已从频率驱动转向并行驱动。GPU 通过大规模并行计算实现了指数级的性能增长，是支撑大语言模型训练和推理的核心硬件。理解 GPU 的工作原理和性能特征，是设计高效 ML 系统的必要条件。

%% ========== GPU 架构 ==========

